% =================================================================
% Bachelorarbeit -- EINE einzige Datei.
%
% Auf Overleaf: nur diese Datei hochladen und uebersetzen.
% Compiler: pdfLaTeX. Kein biber, kein bibtex noetig -- das
% Literaturverzeichnis steht am Ende direkt im Dokument.
%
% Abbildungen kommen erst in den spaeteren Kapiteln dazu. Sobald es
% so weit ist, den Inhalt des Ordners abbildungen/ mit hochladen;
% \graphicspath findet sie dann ueber den zweiten Eintrag.
%
% Gliederung nach BA_V3 (Stand 3. Oktober 2026), die ihrerseits die
% Kapitelreihenfolge aus V2 unveraendert uebernimmt.
% =================================================================
\documentclass[12pt,a4paper]{scrreprt}

\usepackage[utf8]{inputenc}
\usepackage[T1]{fontenc}
\usepackage[ngerman]{babel}
\usepackage{lmodern}
\usepackage{microtype}

\usepackage[a4paper,left=3cm,right=2.5cm,top=2.5cm,bottom=2.5cm]{geometry}
\usepackage{setspace}
\onehalfspacing

\usepackage{graphicx}
\graphicspath{{../abbildungen/}{./}}
\usepackage{booktabs}
\usepackage{amsmath}
\usepackage[hidelinks]{hyperref}

\begin{document}

% -----------------------------------------------------------------
% Titelseite
% -----------------------------------------------------------------
\begin{titlepage}
\centering
\vspace*{1.5cm}

{\Large Ostbayerische Technische Hochschule Regensburg\par}
{\large Fakultät Informatik und Mathematik\par}

\vspace{3cm}
{\Huge\bfseries Multi-Layer-Visualisierung von\\[0.35em]
Tiersichtungsdaten in einer\\[0.35em]
Online-Kartenplattform\par}

\vspace{2.5cm}
{\Large Bachelorarbeit\par}

\vspace{2.5cm}
{\large
vorgelegt von\\[0.6em]
{\bfseries Yassir Sbai}\\[0.4em]
Matrikelnummer: 3341177\par}

\vspace{2cm}
{\large Betreuer und Prüfer: Prof.\ Dr.\ Jan Dünnweber\par}

\vfill
{\large Regensburg, \today\par}
\end{titlepage}

\tableofcontents

% =================================================================
\chapter{Einleitung}
\label{kap:einleitung}

\section{Motivation und Problemstellung}
\label{sec:motivation}

Die Erforschung und der Schutz frei lebender Wildtierpopulationen sind
zunehmend auf digitale Werkzeuge angewiesen, die räumliche und zeitliche
Bewegungsmuster sichtbar und analysierbar machen. Bei großen,
potenziell konfliktträchtigen Arten wie Bären fallen dabei
Datenquellen sehr unterschiedlicher Art an. Besenderte Tiere liefern
GPS-Telemetriedaten, also fortlaufende Positionen eines bekannten
Individuums. Wildkameras liefern Einzelbilder mit Zeitstempel und
Standort, aber ohne Angabe, welches Tier ausgelöst hat.
Simulationsmodelle schließlich erzeugen rechnerisch erzeugte
Bewegungspfade, die weder das eine noch das andere sind. Damit diese
Daten für ein Forschungsteam nutzbar werden, müssen sie nicht nur
gespeichert, sondern verständlich dargestellt werden.

Der Anwendungskontext dieser Arbeit umfasst zwei Regionen. In
\textbf{Japan} haben sich Begegnungen zwischen Menschen und Bären in
den vergangenen Jahren gehäuft. Behörden und Gemeinden reagieren unter
anderem mit technischen Maßnahmen: Wildkameras überwachen bekannte
Wechsel, und akustische oder optische Geräte sollen Bären von
Siedlungen fernhalten~\cite{lorenz2025}. In den \textbf{rumänischen
Karpaten} wurden dagegen im Rahmen einer Telemetriestudie dreizehn
Europäische Braunbären (\textit{Ursus arctos}) über mehrere Jahre mit
Halsbandsendern begleitet~\cite{pop2018}. Beide Bestände betreffen
dieselbe Fragestellung -- wo halten sich die Tiere auf und wie bewegen
sie sich --, liegen aber in völlig verschiedener Form vor.

Im Rahmen einer Kooperation zwischen der Ostbayerischen Technischen
Hochschule Regensburg und einer Partnerhochschule sind im Umfeld dieser
Fragestellung bereits drei eigenständige Teilsysteme entstanden:

\begin{itemize}
  \item ein System zur Verwaltung und Kartierung von
        Bärenabschreckungsmaßnahmen in Japan~\cite{lorenz2025},
  \item ein Modell zur Simulation von Bärenbewegungen auf Basis realer
        GPS-Trajektorien aus den Karpaten~\cite{braeu2025},
  \item ein System zur automatisierten Extraktion von Wildtierbildern
        aus E-Mail-Postfächern~\cite{lelchat2026}.
\end{itemize}

Diese Teilsysteme entstanden in getrennten studentischen Arbeiten. Sie
wurden unabhängig voneinander implementiert, hielten ihre Daten jeweils
eigenständig -- teils in einer MariaDB, teils in lokalen
SQLite-Dateien, teils in CSV-Dateien -- und brachten jeweils eine
eigene Kartenansicht mit, die auf der JavaScript-Bibliothek
Leaflet beziehungsweise deren Python-Anbindung Folium beruhte.

Daraus ergibt sich die eigentliche Problemstellung, und sie betrifft
nicht die Speicherung der Daten, sondern ihre \emph{gemeinsame
Betrachtbarkeit}. Wer beurteilen möchte, ob ein geplanter Standort für
ein Abschreckungsgerät sinnvoll gewählt ist, müsste die
Kamerastandorte, die dort registrierten Sichtungen, die simulierten
Bewegungspfade und die reale Geländebeschaffenheit gleichzeitig vor
Augen haben. Stattdessen muss zwischen mehreren Anwendungen gewechselt
und das Ergebnis im Kopf übereinandergelegt werden.

Hinzu kommt eine zweite, methodische Lücke: Es gab bislang keine
Möglichkeit, ein \emph{Simulationsergebnis gegen Realdaten zu prüfen}.
Die Simulation erzeugte Pfade, aber niemand konnte sehen, ob diese
Pfade denen echter Tiere überhaupt ähneln. Ein Modell, dessen Ausgabe
sich nicht mit der Wirklichkeit vergleichen lässt, ist als
Planungsgrundlage wenig wert.

Dass die Antwort auf beide Punkte eine Karte sein muss und keine
Tabelle, ergibt sich aus der Natur der Fragen. Ob ein simulierter Pfad
an einem Dorf vorbeiführt oder mitten hindurch, ob zwei Kameras
denselben Waldrand abdecken oder zwei verschiedene, ob ein Tier einem
Tal folgt oder es quert -- all das sind räumliche Beziehungen. In einer
Tabelle von Koordinatenpaaren sind sie vorhanden, aber nicht erkennbar.

\section{Zielsetzung der Arbeit}
\label{sec:zielsetzung}

Ziel dieser Arbeit ist die Konzeption und Implementierung einer
einheitlichen Kartenplattform, die die genannten Teilsysteme in einer
gemeinsamen, mehrschichtigen Visualisierung zusammenführt und um
Werkzeuge ergänzt, mit denen sich Bewegungen darstellen und bewerten
lassen. Die Ziele liegen auf zwei Ebenen.

\subsection*{Ziel 1 -- Systemintegration}

Die erste Ebene betrifft die Zusammenführung der bestehenden Systeme zu
einer Anwendung. Dazu gehören vier Teilziele.

\paragraph{Technologiewechsel der Kartenvisualisierung.}
Die bisherige, auf Folium und Leaflet beruhende Kartendarstellung wird
durch MapLibre GL JS abgelöst~\cite{maplibre}. Dabei handelt es sich um
eine quelloffene, WebGL-gestützte Kartenbibliothek, die Vektordaten
erst im Browser zeichnet. Gegenüber dem bisherigen Ansatz verbessert
das die Darstellungsleistung bei großen Punkt- und Pfadmengen, erlaubt
eine einheitliche, deklarative Beschreibung aller Ebenen und eröffnet
zusätzliche Darstellungsformen wie eine Geländedarstellung mit
Höhenrelief.

\paragraph{Systemintegration auf einer gemeinsamen Karte.}
Die drei Teilsysteme -- Bärenabschreckung Japan, Bewegungssimulation
Karpaten und Bildextraktion -- werden in einer Anwendung mit einer
einzigen, gemeinsam genutzten Kartenkomponente zusammengeführt. Die
einzelnen Datenebenen, etwa reale Sichtungen aus Japan, GPS-Pfade aus
Rumänien, Kamerastandorte, Abschreckungsgeräte und simulierte
Bewegungen, lassen sich unabhängig voneinander ein- und ausblenden. Zum
Vergleich werden darüber hinaus externe Datenquellen eingebunden,
namentlich Sichtungsdaten zu Wolf und Luchs sowie Bärenvorkommen aus
der Global Biodiversity Information Facility~\cite{gbif}.

\paragraph{Vereinheitlichung der Datenhaltung.}
Die zuvor teils lokal in SQLite und CSV-Dateien gehaltenen Daten werden
in eine zentrale, gemeinsam nutzbare MariaDB überführt. Das ermöglicht
konsistenten Mehrbenutzerzugriff und serverseitige Filterung, statt
jeden Datensatz vollständig in die Anwendung zu laden.

\paragraph{Interaktions- und Filterkonzept.}
Forschende sollen einzelne Bären, Zeiträume, Jahreszeiten und
Darstellungsarten auswählen können, ohne dass die Anwendung vollständig
neu geladen wird.

\subsection*{Ziel 2 -- Analysewerkzeuge}

Die zweite Ebene ist der inhaltliche Kern der Arbeit. Sie geht über die
Darstellung hinaus und macht die Simulation überprüfbar.

\paragraph{Animierte Darstellung von Bewegungen.}
Bewegungen werden nicht nur als fertige Linie, sondern im zeitlichen
Ablauf dargestellt. Erst dadurch wird erkennbar, wann ein Tier
unterwegs ist und wann es sich kaum bewegt.

\paragraph{Synchrone Darstellung von echter und simulierter Bewegung.}
Zu einer realen, aufgezeichneten Bewegung wird eine simulierte Bewegung
mit identischem Startpunkt und identischem Zeitraster erzeugt und
gleichzeitig animiert. Damit lässt sich unmittelbar beobachten, worin
Modell und Wirklichkeit übereinstimmen und worin nicht. Für diese
Darstellung musste die Simulation von einzelnen Kandidatenschritten auf
durchgehende Bewegungspfade erweitert werden.

\paragraph{Geländebewusste Richtungswahl.}
Die Richtungswahl der Simulation soll nicht nur bewerten, wie das
Gelände am Ende eines Schrittes beschaffen ist, sondern erkennen, in
welcher Richtung zusammenhängendes, für Bären attraktives Gelände
liegt. Dafür wird ein Verfahren übertragen, das ursprünglich zur
Bestimmung von Fußpositionen auf Sensormatten entwickelt
wurde~\cite{geyer2026}. Ergänzend wird die Richtungstreue, also die
Neigung eines Tieres, eine einmal eingeschlagene Richtung
beizubehalten, aus den realen Karpaten-Daten gemessen und in das Modell
übernommen, statt sie zu schätzen.

\paragraph{Quantitativer Vergleich beider Datenbestände.}
Da sich aus den japanischen Kameradaten keine Bewegungsbahnen bilden
lassen, wird ein Vergleich auf einer Ebene gesucht, die beide Bestände
tragen: dem Tagesrhythmus. Über den Überlappungskoeffizienten zweier
Aktivitätsverteilungen~\cite{ridout2009} wird quantifiziert, wie
ähnlich sich die Tiere beider Regionen im Tagesverlauf verhalten.

\paragraph{Nachvollziehbarkeit.}
Durchgehend gilt: Jede in dieser Arbeit berichtete Zahl soll sich mit
einem im Projekt enthaltenen Skript reproduzieren lassen, und die
Simulation soll bei gleicher Zufallszahlenvorgabe stets dasselbe
Ergebnis liefern. Die Vorarbeit kannte keinen festen Startwert für den
Zufallsgenerator; Ergebnisse waren damit nicht wiederholbar.

\subsection*{Abgrenzung}

Einige naheliegende Erwartungen werden ausdrücklich ausgeschlossen.

Es wird \textbf{kein neues Bewegungsmodell von Grund auf entwickelt}.
Das Geschwindigkeits- und Schrittlängenmodell aus der Vorarbeit von
Bräu~\cite{braeu2025} wird unverändert übernommen. Erweitert wird
ausschließlich die Richtungswahl; die Abweichungen von der Vorarbeit
werden einzeln dokumentiert und begründet.

Es wird \textbf{kein Verfahren zur Artbestimmung aus Bildern}
entwickelt. Ob auf einer Kameraaufnahme ein Bär zu sehen ist, wird als
gegeben vorausgesetzt; die Beschaffung der Bilder ist Gegenstand der
Vorarbeit von Lelchat~\cite{lelchat2026}.

Die Anwendung verarbeitet \textbf{keine Echtzeitdaten}. Sie arbeitet
auf abgelegten Beständen; neu eingegangene Bilder werden erkannt, aber
es findet keine Live-Auswertung eines laufenden Kamerastroms statt.

Die Simulation ist \textbf{keine Vorhersage}. Sie erzeugt plausible
Bewegungen unter den Annahmen des Modells, trifft aber keine Aussage
darüber, wo sich ein bestimmtes Tier zu einem bestimmten Zeitpunkt
tatsächlich aufhalten wird. Sie wird entsprechend nicht an
Einzelprognosen gemessen, sondern an statistischen Eigenschaften wie
Tagesstrecke, Netto-Verlagerung und Geländenutzung.

\section{Aufbau der Arbeit}
\label{sec:aufbau}

\textbf{Kapitel~\ref{kap:einleitung}} hat Ausgangslage, Problemstellung
und Ziele umrissen.

\textbf{Kapitel 2} legt die Grundlagen dar, die zum Verständnis der
folgenden Kapitel nötig sind: Geodaten und Web-Kartenvisualisierung,
ein Vergleich gängiger Kartenbibliotheken mit Schwerpunkt auf MapLibre
GL JS, die bewegungsökologischen Begriffe rund um
Step-Selection-Funktionen~\cite{thurfjell2014}, die Methodik der
Kamerafallen-Auswertung sowie die gewichtete Schwerpunktbildung, auf
der die neue Richtungswahl beruht.

\textbf{Kapitel 3} analysiert die bestehenden Systeme: ihre Architektur,
ihre Datenstrukturen und die technischen Grenzen, die sich bei der
Übernahme gezeigt haben. Zusätzlich werden die Eigenschaften der beiden
Datenbestände gemessen und belegt, da mehrere Entwurfsentscheidungen
der Umsetzung erst dadurch begründbar werden.

\textbf{Kapitel 4} beschreibt die Umsetzung: die Kartenvisualisierung
mit MapLibre GL JS, die Integration der Teilsysteme und die
Vereinheitlichung der Datenbank, die Einbindung externer
Vergleichsdaten, die Bewegungsanimation, die Übertragung der Simulation
auf die japanischen Kameradaten einschließlich Geländegewichtung und
Offline-Betrieb, den statistischen Tagesrhythmus-Vergleich sowie die
Interaktions- und Visualisierungskonzepte. Die jeweiligen
Messergebnisse stehen dort, wo auch das zugehörige Verfahren
beschrieben wird.

\textbf{Kapitel 5} bewertet das entwickelte System, legt seine
Limitationen offen und ordnet die Messergebnisse wissenschaftlich ein.
Dazu gehört ausdrücklich auch ein negativer Befund zur Wirksamkeit des
übertragenen Verfahrens.

\textbf{Kapitel 6} fasst die Ergebnisse zusammen und gibt einen Ausblick
auf mögliche Erweiterungen.

% =================================================================
% Weitere Kapitel werden hier unten angehaengt. Struktur nach V3:
%
% \chapter{Grundlagen}
%   2.1 Geodaten und Web-Kartenvisualisierung
%   2.2 Web-Kartenbibliotheken (insb. MapLibre GL JS)
%   2.3 Bewegungsoekologie und Step-Selection-Funktionen
%   2.4 Kamerafallen-Methodik und Aktivitaetsmuster-Analyse
%   2.5 Gewichtete Schwerpunktbildung und Instanztrennung
%
% \chapter{Analyse der bestehenden Systeme}
%   3.1 Systemuebersicht und Komponenten
%   3.2 Simulationsdaten
%   3.3 Datengrundlage und ihre gemessenen Eigenschaften
%
% \chapter{Implementierung}
%   4.1 Kartenvisualisierung mit MapLibre GL JS
%   4.2 Integration der Teilsysteme
%   4.3 Externe Vergleichsdatenquellen (Wolf/Luchs, GBIF)
%   4.4 Bewegungsanimation
%   4.5 Uebertragung auf japanische Kameradaten
%   4.6 Statistischer Tagesrhythmus-Vergleich
%   4.7 Interaktions- und Visualisierungskonzepte
%   4.8 Robustheit und Betriebsfaehigkeit
%
% \chapter{Diskussion}
%   5.1 Bewertung des entwickelten Systems
%   5.2 Limitationen
%   5.3 Wissenschaftliche Einordnung des Delta-Ergebnisses
%   5.4 Wirksamkeit der uebertragenen Methode
%
% \chapter{Fazit und Ausblick}
% =================================================================

% =================================================================
% Literaturverzeichnis
%
% Noch zu klaeren:
%  - Art der Arbeit von Lelchat beim Betreuer erfragen
%  - finale Angaben zum Journal-Artikel, sobald erschienen
%  - Seitenzahlen der Konferenzfassung
%  - Ridout und Linkie gegen das Original pruefen
%  - Primaerquelle fuer die Fallzahlen in Japan ergaenzen
%  - genaue GBIF-Download-DOI und BfN-Quelle ergaenzen
% =================================================================
\begin{thebibliography}{99}

\bibitem{braeu2025}
F.~Bräu,
\emph{Analyse und Simulation von Bärenbewegungen}.
Bachelorarbeit, Ostbayerische Technische Hochschule Regensburg,
Fakultät Informatik und Mathematik, August 2025.

\bibitem{lorenz2025}
M.~Lorenz,
\emph{Entwicklung und Implementierung eines interaktiven GIS-Systems zur
Kartierung von Wildtierdaten mit Fokus auf
Bärenabschreckungsmaßnahmen in Japan}.
Bachelorarbeit, Ostbayerische Technische Hochschule Regensburg,
Fakultät Informatik/Mathematik, Juni 2025.

\bibitem{lelchat2026}
A.~Lelchat und J.~Dünnweber,
\emph{Extracting wildlife pictures from emails}.
Ostbayerische Technische Hochschule Regensburg, 2026.

\bibitem{geyer2026}
J.~Geyer, M.~Melzer, J.~Dünnweber und A.~Sarkar,
"`Human Step Movement Imitation for Training Privacy-Aware Intelligent
Environments"', in \emph{2026 22nd International Conference on
Intelligent Environments (IE)}, Juni 2026.
DOI: 10.1109/IE69249.2026.11539056.

\bibitem{duennweber2026}
J.~Dünnweber und J.~Geyer,
"`Walking in Every Direction: Aperiodic Training Data and Anisotropic
Kernels for Step Movement Analysis"',
\emph{Journal of Ambient Intelligence and Smart Environments},
im Erscheinen.

\bibitem{pop2018}
I.~M.~Pop, L.~Bereczky, S.~Chiriac, R.~Iosif, A.~Nita, V.~D.~Popescu
und L.~Rozylowicz,
"`Movement ecology of brown bears (\emph{Ursus arctos}) in the Romanian
Eastern Carpathians"', \emph{Nature Conservation}, Bd.~26,
S.~15--31, 2018. DOI: 10.3897/natureconservation.26.22955.

\bibitem{thurfjell2014}
H.~Thurfjell, S.~Ciuti und M.~S.~Boyce,
"`Applications of step-selection functions in ecology and
conservation"', \emph{Movement Ecology}, Bd.~2, Nr.~4, 2014.
DOI: 10.1186/2051-3933-2-4.

\bibitem{ridout2009}
M.~S.~Ridout und M.~Linkie,
"`Estimating overlap of daily activity patterns from camera trap
data"', \emph{Journal of Agricultural, Biological, and Environmental
Statistics}, Bd.~14, Nr.~3, S.~322--337, 2009.

\bibitem{maplibre}
MapLibre-Projekt,
\emph{MapLibre GL JS}.
\url{https://maplibre.org/maplibre-gl-js/docs/}

\bibitem{gbif}
Global Biodiversity Information Facility,
\emph{GBIF Occurrence Data}.
\url{https://www.gbif.org}

\bibitem{openstreetmap}
OpenStreetMap-Mitwirkende,
\emph{OpenStreetMap}. Daten unter ODbL.
\url{https://www.openstreetmap.org}

\bibitem{overpass}
OpenStreetMap Wiki,
\emph{Overpass API}.
\url{https://wiki.openstreetmap.org/wiki/Overpass_API}

\bibitem{epsg3844}
IOGP Geomatics Committee,
\emph{EPSG Geodetic Parameter Registry}, Code 3844
(Pulkovo 1942(58) / Stereo70).
\url{https://epsg.org/}

\end{thebibliography}

\end{document}
